% Einheit 3 von 5 – Prisma (bank/koerper/e3.jsonl, zusammenbau v0.1)
%% TODO Zweigzeile Teil 1: was hier gelernt wird (Satz in Schülersprache aus dem Einheitstitel) – Einheit 3
\einheitenkopf[e3]{Einheit 3 von 5 · Prisma}
\zweigzeile{ab Kl. 8 (am Gymnasium ab Kl. 7) · P10}
\setcounter{aufgabe}{19}

%% TODO Titel: Ich-kann-Satz fehlt in der Bank (Platzhalter: Kettenname) – Nr. 20
%% TODO Anweisung: ein Satz über der ersten Teilaufgabe fehlt in der Bank – Nr. 20
\begin{aufgabe}{Prisma}
%% TODO \rechenplatz{n}: Schrittzahl des Grundfalls fehlt in der Bank (nur bei fester Schreibform, 2.3 b)
\begin{teile}
\teil Ein Zelt hat die Form dieses Prismas. Färbe die Grundfläche und zeichne die Höhe des Prismas farbig nach.

\prismadreieck{4}{2.5}{2}{}{}{}
\teil Prisma: Grundfläche $18$ cm², Höhe $7$ cm – Volumen? \leerfeld[cm³]
\teil Prisma mit rechteckiger Grundfläche $7$ cm × $3$ cm, Höhe $5$ cm – Volumen? Grundfläche \leerfeld[cm²] , Volumen \leerfeld[cm³]
\teil Dreiecksprisma: Dreieck mit Grundseite $8$ cm und Höhe $5$ cm; Höhe des Prismas $9$ cm – Volumen? Grundfläche \leerfeld[cm²] , Volumen \leerfeld[cm³]
\teil Prisma mit Trapez als Grundfläche: parallele Seiten $7$ cm und $3$ cm, Abstand $4$ cm; Höhe des Prismas $6$ cm – Volumen? Grundfläche \leerfeld[cm²] , Volumen \leerfeld[cm³]
\teil Ein Zelt liegt wie im Bild: vorn ein Dreieck, $2{,}4$ m breit und $1{,}6$ m hoch; das Zelt ist $3$ m lang. Wie viel Luft ist im Zelt?

\prismadreieck{4}{2.5}{2}{$2{,}4$ m}{$1{,}6$ m}{$3$ m}
\leerfeld[m³]
\teil Die Grundfläche eines Prismas ist ein Dreieck mit den Seiten $4$, $7$ und $8$ cm, das Prisma ist $5$ cm hoch – Mantel? \leerfeld[cm²]
\teil Dreiecksprisma: Das Dreieck hat einen rechten Winkel zwischen den Seiten $3$ cm und $4$ cm, die lange Seite ist $5$ cm; das Prisma ist $7$ cm hoch – Oberfläche? \leerfeld[cm²]
\teil Vom Netz eines Dreiecksprismas sind drei Rechtecke nebeneinander gezeichnet: $4$ cm × $7$ cm, $5$ cm × $7$ cm und $6$ cm × $7$ cm. Zeichne das Netz ab und ergänze, was fehlt.

\rechenplatz{5}
\teil Prisma: Volumen $210$ cm³, Grundfläche $35$ cm² – Höhe? \leerfeld[cm]
\teil Ein Brunnentrog hat die Form eines Prismas; Grund- und Deckfläche sind gleichseitige Dreiecke mit der Seite $1{,}80$ m und der Höhe $1{,}56$ m. Das Prisma ist $0{,}75$ m hoch. Berechne den Flächeninhalt der Deckfläche und weise nach, dass das Volumen etwa $1{,}1$ m³ beträgt \hfill (P10 2019 OS). Deckfläche \leerfeld[m²] , Volumen \leerfeld[m³]
\end{teile}
\end{aufgabe}

%% TODO Titel: Ich-kann-Satz fehlt in der Bank (Platzhalter: Kettenname) – Nr. 21
%% TODO Anweisung: ein Satz über der ersten Teilaufgabe fehlt in der Bank – Nr. 21
\begin{aufgabe}{Prisma – Fehler finden, Begründen, Anwendung, Darstellungswechsel}
\begin{teile}
\teil Dreiecksprisma: Dreieck mit Grundseite $7$ cm und Höhe $6$ cm, das Prisma ist $5$ cm hoch. Emil rechnet: \rechnung{7 \cdot 6 &= 42 \text{ cm}^2 \\ 42 \cdot 5 &= 210 \text{ cm}^3} Wo ist der Fehler?
\teil Warum rechnet man beim Prisma Grundfläche mal Höhe? Denk an einen Stapel.
\teil Ein Dachboden hat die Form eines liegenden Dreiecksprismas: $9$ m breit, $3{,}5$ m hoch und $11$ m lang. Ein Heizgerät reicht für Räume bis $180$ m³. Reicht es für den ausgebauten Dachboden?
\teil Das Dreieck ist gleichschenklig, die Schenkel sind je $5$ cm lang. Skizziere das Netz des Prismas mit allen Maßen.

\prismadreieck{4}{1.5}{2.5}{$8$ cm}{$3$ cm}{$7$ cm} \rechenplatz[halb]{4}
\end{teile}
\end{aufgabe}
